% Generated from locale/tr/content/first-order-logic/syntax-and-semantics/free-vars-sentences.tex; do not hand-edit.


\olfileid[tr]{fol}{syn}{fvs}

\olsection{Serbest \printtoken{P}{variable} ve \printtoken{P}{sentence}}

\begin{defn}[Bir !!a{variable} için serbest oluşumlar]
\ollabel{defn:free-occ}
Bir !!a{variable} için bir !!a{formula} içindeki \emph{serbest}
oluşumlar aşağıdaki gibi tümevarımsal olarak tanımlanır:
\begin{enumerate}
\item \indcase*{!A}{$!A$ atomiktir}{$\indfrm$ içindeki bütün
  !!{variable} oluşumları serbesttir.}

\tagitem{prvNot}{\indcase{!A}{\lnot !B}{$\indfrm$ içindeki serbest
  !!{variable} oluşumları tam olarak $!B$'dekilerdir.}}{}

\item \indcase{!A}{(!B \ast !C)}{$\indfrm$ içindeki serbest
  !!{variable} oluşumları, $!B$'dekiler ile~$!C$'dekilerin birlikte
  oluşturduğu kümedir.}

\tagitem{prvAll}{\indcase{!A}{\lforall[x][!B]}{$\indfrm$ içindeki serbest
  !!{variable} oluşumları, $x$ oluşumları dışında $!B$'deki oluşumların
  tümüdür.}}{}

\tagitem{prvEx}{\indcase{!A}{\lexists[x][!B]}{$\indfrm$ içindeki serbest
  !!{variable} oluşumları, $x$ oluşumları dışında $!B$'deki oluşumların
  tümüdür.}}{}
\end{enumerate}
\end{defn}

\begin{defn}[Bağlı Değişkenler]
Bir $!A$ formülündeki bir !!a{variable} oluşumu serbest değilse
\emph{bağlı}dır.
\end{defn}

\begin{prob}
\olref[fol][syn][fvs]{defn:free-occ} doğrultusunda bağlı değişken
oluşumlarının tümevarımsal bir tanımını verin.
\end{prob}

\begin{defn}[Kapsam]
\iftag{prvAll}{$\lforall[x][!B]$, bir $!A$ formülü içinde bir alt formül
  oluşumuysa, $!A$ içindeki buna karşılık gelen $!B$ oluşumuna
  $\lforall[x]$'in karşılık gelen oluşumunun \emph{kapsamı} denir.
  \iftag{prvEx}{$\lexists[x]$ için de benzer biçimde.}{}}{$\lexists[x][!B]$,
  bir $!A$ formülü içinde bir alt formül oluşumuysa, $!A$ içindeki buna
  karşılık gelen $!B$ oluşumuna $\lexists[x]$'in karşılık gelen
  oluşumunun \emph{kapsamı} denir.}

$!B$, $!A$ içinde bir
\iftag{prvAll}{$\lforall[x]$\iftag{prvEx}{ ya da
    $\lexists[x]$}{}}{$\lexists[x]$} niceleyici oluşumunun kapsamıysa,
$!B$ içindeki serbest $x$ oluşumları
\iftag{prvAll}{$\lforall[x][!B]$\iftag{prvEx}{ ve
$\lexists[x][!B]$}{}}{$\lexists[x][!B]$} içinde bağlıdır. Bu oluşumlara,
söz konusu niceleyici oluşumu \emph{tarafından bağlanmış} deriz.
\end{defn}

\begin{ex}
Şu formülü ele alalım:
\[
\lexists[\Obj v_0][\underbrace{\Atom{\Obj A^2_0}{\Obj v_0,\Obj v_1}}_{!B}] 
\]
$!B$, $\lexists[\Obj v_0]$'ın kapsamını gösterir. Niceleyici, $!B$
içindeki $\Obj v_0$ oluşumunu bağlar; fakat $\Obj v_1$ oluşumunu
bağlamaz. Dolayısıyla $\Obj v_1$ bu durumda serbest bir değişkendir.


Şimdi bunun daha karmaşık bir !!{formula}~$!A$ içinde nasıl
işleyebileceğini görebiliriz:
\[
\lforall[\Obj v_0][\underbrace{(\Atom{\Obj A^1_0}{\Obj v_0} \lif
    \Atom{\Obj A^2_0}{\Obj v_0, \Obj v_1})}_{!B}] \lif \lexists[\Obj
  v_1][\underbrace{(\Atom{\Obj A^2_1}{\Obj v_0, \Obj v_1} \lor \lforall[\Obj v_0][\overbrace{\lnot \Atom{\Obj A^1_1}{\Obj v_0}}^{!D}])}_{!C}]
\]
$!B$, ilk $\lforall[\Obj v_0]$'ın; $!C$, $\lexists[\Obj v_1]$'in; $!D$
ise ikinci $\lforall[\Obj v_0]$'ın kapsamıdır. İlk
$\lforall[\Obj v_0]$, $!B$ içindeki $\Obj v_0$ oluşumlarını;
$\lexists[\Obj v_1]$, $!C$ içindeki $\Obj v_1$ oluşumunu; ikinci
$\lforall[\Obj v_0]$ ise $!D$ içindeki $\Obj v_0$ oluşumunu bağlar.
$\Obj v_1$'in ilk oluşumu ile $\Obj v_0$'ın dördüncü oluşumu $!A$
içinde serbesttir. $\Obj v_0$'ın son oluşumu $!D$ içinde serbest, fakat
$!C$ ve~$!A$ içinde bağlıdır.
\end{ex}

\begin{defn}[Tümce]
!!^a{formula}~$!A$, \article{sentence} \emph{!!{sentence}} olarak
adlandırılır ancak ve ancak !!{variable}s{}in hiçbir serbest oluşumunu içermiyorsa.
\end{defn}

% add examples!

